\paragraph{Passage terminology.}
\label{tab:terminology}
A central height is \emph{admissible} if it belongs to \(\cH_p\);
the associated half-orbits reach the windows \(W_\pm\).
A \emph{captured} periodic orbit is one satisfying the Hausdorff
condition of \cref{prop:isolation-capture-new}. Its central height
is admissible by \cref{prop:central-displacement-interval-new}.

A \emph{prefix} runs from the central point to a prescribed radius,
and a \emph{tail} from an incoming height to \(z_*\).
The \emph{projected state} is the boundary-flow coordinate \(w\).
A half-orbit is \emph{thin} while its stated normalized height is
below the chosen threshold; its \emph{thin exit} is the first outward
crossing of that threshold. The subsequent \emph{landing} is the first
hit of \(z=m_{\rm land}|\eta|^2\).

A \emph{weak root} approaches the central line relative to the current
coefficient scale. The \emph{fold} is a double root of the positive
secondary drift; a small positive minimum on its no-root side is the
\emph{bottleneck}. The \emph{far root} is the negative branch's first
primary root \(X_*=1+O(r+s)\) in
\cref{sec:negative-branch-far-root}.
